\documentclass[11pt,a4paper]{article}
\usepackage{amsmath,amssymb,amsthm}
\usepackage{hyperref}
\usepackage{geometry}
\usepackage{booktabs}
\usepackage{mdframed}
\usepackage{xcolor}
\usepackage{draftwatermark}
\SetWatermarkText{DRAFT}
\SetWatermarkScale{3}
\SetWatermarkColor[gray]{0.92}
\geometry{margin=2.5cm}

\newtheorem{theorem}{Theorem}
\newtheorem{lemma}[theorem]{Lemma}
\newtheorem{definition}[theorem]{Definition}
\newtheorem{remark}[theorem]{Remark}
\newtheorem{axiom*}[theorem]{Axiom (externally assumed)}

\title{\textbf{DRAFT --- Under Review and Validation}\\[0.4em]
A Lean~4 Scaffolding for Riemann Hypothesis Equivalences:\\
\large Formal Criterion Infrastructure and Bridge Lemma Discovery}
\author{Xavier Fresquet \and Gérard Biau\\[0.3em]
\normalsize Sorbonne Université, Paris \& Abu Dhabi\\
\normalsize \textit{Under review and validation --- not definitive results}}
\date{June 2026 --- Preliminary version}

\begin{document}
\maketitle

\begin{mdframed}[backgroundcolor=yellow!15, linecolor=orange!60, linewidth=1pt]
\textbf{Notice.} This is a \textbf{draft} report.  Results are preliminary
and under active review and validation by Biau \& Fresquet.  Nothing here
should be read as a definitive mathematical claim.
\end{mdframed}

\vspace{1em}

\begin{abstract}
We describe a Lean~4 project that formalizes several classical equivalences of
the Riemann Hypothesis (RH) and builds the infrastructure needed to search for
bridge lemmas between them.  The project does \emph{not} prove RH, nor does it
claim to be close to doing so.  Its aim is more modest: to provide a rigorous
Lean environment in which (i) each equivalent formulation is stated with correct
types, (ii) the logical chain from a certificate structure to RH is
machine-checked, and (iii) the precise ``sorry'' gaps are recorded as honest
targets for future proof search.
\end{abstract}

\section{Background}

The Riemann Hypothesis asserts that every non-trivial zero $s$ of the Riemann
zeta function $\zeta(s)$ satisfies $\mathrm{Re}(s) = \tfrac{1}{2}$.  Several
classical criteria are known to be equivalent to RH:

\begin{itemize}
  \item \textbf{Nyman--Beurling--Báez-Duarte}: RH holds if and only if the
        characteristic function $\chi_{(0,1)}$ belongs to the closure in
        $L^2(0,1)$ of the span of the functions
        $\rho_n(x) = \{n/x\}$ (fractional part), $n \geq 1$.
  \item \textbf{Li's positivity criterion}: RH holds if and only if all Li
        coefficients $\lambda_n \geq 0$.
  \item \textbf{Weil explicit formula and positivity}: RH is equivalent to a
        positivity condition on a certain Weil kernel functional.
  \item \textbf{Hilbert--Pólya spectral criterion}: RH would follow from the
        existence of a self-adjoint operator whose spectrum is the sequence of
        imaginary parts of the non-trivial zeros.
\end{itemize}

\section{Lean Architecture}

The project uses Lean~4 (v4.30.0) with Mathlib and is structured as a Lake
library under \texttt{riemann/}:

\begin{itemize}
  \item \texttt{Basic/Zeta.lean} --- definitions of $\zeta(s)$, $\xi(s)$, and
        the critical strip; bridged to Mathlib's \texttt{riemannZeta}.
  \item \texttt{Basic/CriticalStrip.lean} --- the formal statement of RH as a
        Lean \texttt{Prop}, zero symmetry, and the zero-free region near
        $\mathrm{Re}(s) = 1$.
  \item \texttt{Criteria/NymanBeurling.lean} --- the Nyman--Beurling--Báez-Duarte
        criterion; certificate structures; the main bridge theorem
        (sorry-free conditional chain).
  \item \texttt{Criteria/Li.lean}, \texttt{Weil.lean}, \texttt{Spectral.lean}
        --- Li, Weil, and Hilbert--Pólya criterion stubs.
  \item \texttt{Finite/GramMatrices.lean}, \texttt{TailBounds.lean} ---
        finite Gram certificates and transfer theorems.
  \item \texttt{Discovery/BridgeLemmas.lean} --- candidate bridge lemmas as
        sorry stubs for Aristotle.
  \item \texttt{Discovery/FailedAttempts.lean} --- documented proof attempts
        that did not work; a negative-result register.
\end{itemize}

\section{What Is Proved (Sorry-Free)}

\subsection{Formal Criterion Statement}

\begin{definition}[Riemann Hypothesis in Lean]
\[
  \texttt{RiemannHypothesis} \;:=\;
  \forall s : \mathbb{C},\;
  \texttt{isNontrivialZero}\;s \;\to\; \texttt{onCriticalLine}\;s.
\]
\end{definition}

\subsection{Nyman--Beurling Certificate Chain}

The key sorry-free contribution is a complete logical chain from a
\emph{verified approximation rate} to RH, contingent on one externally assumed
equivalence axiom.

\begin{theorem}[Certificate implies Abstract Criterion (sorry-free)]
  Let $v$ be a \texttt{VerifiedApproximationRate}: a sequence $r : \mathbb{N}
  \to \mathbb{R}$ tending to $0$ together with, for every $N$, a certificate
  that the approximation distance satisfies $\mathit{ApproximationDistance}(N)
  \leq r(N)$.  Then the \emph{AbstractCriterion} holds (there exists a rate
  function bounding the approximation distance that tends to zero).
\end{theorem}

\begin{theorem}[Main Bridge (sorry-free modulo one axiom)]
  \label{thm:rh_bridge}
  A \texttt{VerifiedApproximationRate} implies the Riemann Hypothesis.
\end{theorem}

Theorem~\ref{thm:rh_bridge} chains through:
\[
  \texttt{VerifiedApproximationRate}
  \;\xrightarrow{\text{sorry-free}}\;
  \texttt{AbstractCriterion}
  \;\xrightarrow{\text{axiom}}\;
  \texttt{RiemannHypothesis}.
\]

The axiom \texttt{nyman\_beurling\_implies\_rh} encodes the classical analytic
equivalence.  It is quarantined and labeled; it is not presented as a Lean
proof.

\section{What Is Not Proved}

\begin{remark}[ApproximationDistance is undefined]
  The term \texttt{ApproximationDistance} is currently a sorry stub: the
  actual $L^2$ distance from $\chi_{(0,1)}$ to the $N$-dimensional span in
  $L^2(0,1)$ has not been formalized.  This is the next concrete target for
  Mathlib integration (fractional parts, $L^2$ completions).
\end{remark}

\begin{remark}[Analytic criteria are axioms]
  All four RH-equivalent criteria are either sorry stubs (Li, Weil, Spectral)
  or backed by a single named axiom (Nyman--Beurling).  None of the
  equivalences have been proved inside Lean; they are imported from the
  mathematical literature as trusted external facts.
\end{remark}

\begin{remark}[Zero symmetry and zero-free region]
  \texttt{zero\_symmetry} and \texttt{zero\_free\_near\_one} are sorry stubs;
  the Mathlib \texttt{riemannZeta} API needs additional bridging lemmas to
  discharge these.
\end{remark}

\section{Numerical Support}

Three Python experiments provide numerical grounding:
\begin{itemize}
  \item \texttt{gram\_matrix\_search.py} --- computes Gram matrices for the
        Nyman basis and searches for positive-semidefinite certificates.
  \item \texttt{li\_coefficients.py} --- numerical computation and sign check
        of Li coefficients $\lambda_n$ for $n \leq 100$.
  \item \texttt{weil\_kernel\_search.py} --- candidate Weil kernel
        constructions; numerical positivity checks.
\end{itemize}
None of these produce Lean proofs directly; they guide which sorry stubs are
most promising to attack next.

\section{Roadmap}

Future development is planned along the following lines:

\begin{enumerate}
  \item \textbf{Formalize \texttt{ApproximationDistance}}: connect Lean's
        \texttt{Int.fract} and Mathlib's $L^2$ spaces to define the Nyman
        basis and compute the distance formally.
  \item \textbf{Gram matrix certificates}: for small $N$, produce and verify
        positive-semidefinite Gram matrices as Lean \texttt{decide}-level
        certificates, giving a finite approximation guarantee.
  \item \textbf{Tail bound transfer}: formalize the theorem that a sequence of
        finite approximation guarantees, plus a verified tail bound, implies
        the abstract criterion without any axiom.
  \item \textbf{Li criterion infrastructure}: formalize Li's positivity
        definition and bridge it to Mathlib's \texttt{riemannZeta} using the
        functional equation.
  \item \textbf{Aristotle bridge-lemma search}: send the sorry stubs in
        \texttt{Discovery/BridgeLemmas.lean} to Aristotle one at a time;
        record successes and failures in \texttt{FailedAttempts.lean}.
  \item \textbf{Lean 4 / Mathlib upgrade path}: track Mathlib's growing
        complex-analysis library; rebind stubs as the relevant lemmas become
        available.
\end{enumerate}

\section{Next Update}

The following items are planned for the next version of this report, once the
experiment is complete:

\begin{enumerate}
  \item Replace the \texttt{ApproximationDistance} sorry stub with a proper
        Lean definition using Mathlib's fractional-part and $L^2$ APIs.
  \item Present Gram matrix certificates for $N = 5, 10, 20$ verified at the
        Lean kernel level.
  \item Report whether the tail-bound transfer theorem is proved or remains a
        target; include its axiom footprint.
  \item Include the Aristotle log: which stubs were filled, which failed, and
        the tactic where each failure occurred.
  \item Report progress on connecting the Li criterion to Mathlib's
        \texttt{riemannZeta} functional equation.
\end{enumerate}

\section{Conclusion}

This Lean~4 project demonstrates that a rigorous formal scaffold for
RH-equivalent criteria can be built with modest effort, even without a proof
of RH itself.  The value is not proximity to a proof but \emph{precision}: the
exact sorry gaps record, in machine-checkable form, what the mathematical
community currently does not know how to prove inside a formal system.  The
sorry-free bridge theorem (Theorem~\ref{thm:rh_bridge}) shows that the logical
chain from a certificate structure to RH type-checks cleanly, waiting only for
the certificate construction itself.

\end{document}
